\documentclass{article}
\usepackage[utf8]{inputenc}
\usepackage{amsmath,amssymb}
\usepackage[most]{tcolorbox}
\usepackage{geometry}
\geometry{margin=2.5cm}

\newcommand{\step}[1]{\par\noindent\hangindent=3.5em\hangafter=1%
  \makebox[3.5em][l]{\textbf{#1.}}}
\newcommand{\steptag}[1]{\hfill{\small\textsf{[#1]}}}

\newtcolorbox{proofbox}[1]{
  colback=gray!5, colframe=gray!60,
  fonttitle=\bfseries, title={#1},
  breakable, enhanced,
  left=4pt, right=4pt, top=4pt, bottom=4pt,
  before skip=10pt, after skip=10pt
}

\begin{document}

\begin{proofbox}{Cross-pivot cancellation: let P be a rank-3 exact signed ...}
\step{1} Cross-pivot cancellation: let P be a rank-3 exact signed idempotent (square real matrix with P\textasciicircum{}2 = P and all row sums equal to 1), let U = (u\_0,u\_1,u\_2) be an actual-row chart whose rows p\_\{u\_0\}, p\_\{u\_1\}, p\_\{u\_2\} form a basis of the row space, and define coordinates a\_q(i) by p\_i = sum\_q a\_q(i)p\_\{u\_q\} and beta\_r(i) = P\_\{u\_r i\}; then for every pair of distinct indices r, s in \{0,1,2\}, sum\_i beta\_r(i)*a\_s(i) = 0, the sum running over all row indices i of P. \steptag{validated} \\[6pt]
\step{1.1} For every row index k of P, the row p\_k is reproduced by the rows of P with coefficients from row k: p\_k = sum\_i P\_\{k i\} p\_i. Indeed, the j-th coordinate of the right hand side is sum\_i P\_\{k i\} P\_\{i j\} = (P\textasciicircum{}2)\_\{k j\} = P\_\{k j\}, using def-signed-idempotent (P\textasciicircum{}2 = P). \steptag{validated} \\[4pt]
\step{1.2} Specializing the previous row-reproduction identity to k = u\_r and using beta\_r(i) = P\_\{u\_r i\} gives p\_\{u\_r\} = sum\_i beta\_r(i) p\_i. \steptag{validated} \\[4pt]
\step{1.3} Substituting the chart-coordinate expansion p\_i = sum\_q a\_q(i) p\_\{u\_q\} into p\_\{u\_r\} = sum\_i beta\_r(i) p\_i and collecting the finite sums gives p\_\{u\_r\} = sum\_q (sum\_i beta\_r(i)*a\_q(i)) p\_\{u\_q\}. \steptag{validated} \\[6pt]
\step{1.3.1} By the coordinate definition in the root statement, for every row index i one has p\_i = sum\_q a\_q(i) p\_\{u\_q\}, where q ranges over \{0,1,2\}. \steptag{validated} \\[4pt]
\step{1.3.2} Under the root hypotheses, p\_\{u\_r\} = sum\_i beta\_r(i) p\_i: by P\textasciicircum{}2 = P, the j-th coordinate of sum\_i P\_\{u\_r i\} p\_i is sum\_i P\_\{u\_r i\} P\_\{i j\} = (P\textasciicircum{}2)\_\{u\_r j\} = P\_\{u\_r j\}, and beta\_r(i) = P\_\{u\_r i\}. By the root definition of the chart coordinates, each row has p\_i = sum\_q a\_q(i) p\_\{u\_q\}, with q in \{0,1,2\}. Substituting these expansions gives p\_\{u\_r\} = sum\_i beta\_r(i) (sum\_q a\_q(i) p\_\{u\_q\}); since the row-index set is finite and q ranges over the finite set \{0,1,2\}, scalar distributivity and interchange of finite sums give p\_\{u\_r\} = sum\_q (sum\_i beta\_r(i)*a\_q(i)) p\_\{u\_q\}. \steptag{validated} \\[6pt]
\step{1.3.2.1} The premises used in this node are now explicit: node 1.2 supplies p\_\{u\_r\} = sum\_i beta\_r(i) p\_i, and validated node 1.3.1 supplies, for every row index i, the chart expansion p\_i = sum\_q a\_q(i) p\_\{u\_q\} with q in \{0,1,2\}. \steptag{validated} \\[4pt]
\step{1.3.2.2} Substituting the expansion from node 1.3.1 into the equality from node 1.2 gives p\_\{u\_r\} = sum\_i beta\_r(i) (sum\_q a\_q(i) p\_\{u\_q\}). This is just replacement of each p\_i by its defining three-term chart expansion inside the finite row-index sum. \steptag{validated} \\[4pt]
\step{1.3.2.3} For each row index i, scalar distribution gives beta\_r(i) (sum\_q a\_q(i) p\_\{u\_q\}) = sum\_q (beta\_r(i)*a\_q(i)) p\_\{u\_q\}. Hence p\_\{u\_r\} = sum\_i sum\_q (beta\_r(i)*a\_q(i)) p\_\{u\_q\}. \steptag{validated} \\[4pt]
\step{1.3.2.4} The row-index set is finite because P is a square real matrix, and q ranges over the finite set \{0,1,2\}. Therefore the double finite sum may be interchanged and regrouped by q, yielding p\_\{u\_r\} = sum\_q (sum\_i beta\_r(i)*a\_q(i)) p\_\{u\_q\}, which is the claimed collected expansion. \steptag{validated} \\[4pt]
\step{1.4} Because p\_\{u\_0\}, p\_\{u\_1\}, p\_\{u\_2\} form a basis of the row space, coordinates in that chart are unique. The coordinate vector of p\_\{u\_r\} in this ordered basis is the standard vector with q-coordinate delta\_\{q r\}; comparing with the previous expansion yields sum\_i beta\_r(i)*a\_q(i) = delta\_\{q r\} for each q in \{0,1,2\}. \steptag{validated} \\[6pt]
\step{1.4.1} The ordered rows p\_\{u\_0\}, p\_\{u\_1\}, p\_\{u\_2\} are assumed in the root statement to form a basis of the row space; hence any vector in the row space has at most one expansion as sum\_q c\_q p\_\{u\_q\}. \steptag{validated} \\[4pt]
\step{1.4.2} For the chart row p\_\{u\_r\}, the standard expansion is p\_\{u\_r\} = sum\_q delta\_\{q r\} p\_\{u\_q\}, because exactly the r-th coefficient is 1 and the other two coefficients are 0. \steptag{validated} \\[4pt]
\step{1.4.3} Using the row-reproduction specialization p\_\{u\_r\} = sum\_i beta\_r(i)*p\_i and the defining chart-coordinate expansions p\_i = sum\_q a\_q(i)*p\_\{u\_q\}, finite distributivity gives p\_\{u\_r\} = sum\_q (sum\_i beta\_r(i)*a\_q(i))*p\_\{u\_q\}. Comparing this expansion with the standard expansion p\_\{u\_r\} = sum\_q delta\_\{q r\}*p\_\{u\_q\} via uniqueness in the ordered basis p\_\{u\_0\}, p\_\{u\_1\}, p\_\{u\_2\}, each coefficient agrees, so sum\_i beta\_r(i)*a\_q(i) = delta\_\{q r\} for every q. \steptag{validated} \\[6pt]
\step{1.4.3.1} By node 1.2, p\_\{u\_r\} = sum\_i beta\_r(i)*p\_i. By the chart-coordinate definition recorded in node 1.3.1, each p\_i = sum\_q a\_q(i)*p\_\{u\_q\}. Substitution gives p\_\{u\_r\} = sum\_i beta\_r(i)*(sum\_q a\_q(i)*p\_\{u\_q\}); since the row-index set and \{0,1,2\} are finite, bilinearity of scalar multiplication and addition permits distributing beta\_r(i) and interchanging the finite sums, yielding p\_\{u\_r\} = sum\_q (sum\_i beta\_r(i)*a\_q(i))*p\_\{u\_q\}. \steptag{validated} \\[4pt]
\step{1.4.3.2} By node 1.4.2, the same vector p\_\{u\_r\} has the standard basis expansion p\_\{u\_r\} = sum\_q delta\_\{q r\}*p\_\{u\_q\}: the coefficient of p\_\{u\_r\} is 1 and the coefficients of the other two chart rows are 0. \steptag{validated} \\[4pt]
\step{1.4.3.3} Nodes 1.4.3.1 and 1.4.3.2 are two expansions of the same vector p\_\{u\_r\} in the ordered basis p\_\{u\_0\}, p\_\{u\_1\}, p\_\{u\_2\}. By the uniqueness statement in node 1.4.1, the coefficient tuples must be equal coordinatewise. Therefore, for every q in \{0,1,2\}, sum\_i beta\_r(i)*a\_q(i) = delta\_\{q r\}. \steptag{validated} \\[4pt]
\step{1.5} Taking q = s in the coordinate identity gives sum\_i beta\_r(i)*a\_s(i) = delta\_\{s r\}; since r and s are distinct, delta\_\{s r\} = 0, which is exactly the claimed cross-pivot cancellation. \steptag{validated} \\[6pt]
\step{1.5.1} Fix distinct r,s in \{0,1,2\} under the root hypotheses. Since P\textasciicircum{}2 = P, for each coordinate j the j-th coordinate of sum\_i P\_\{u\_r i\} p\_i is sum\_i P\_\{u\_r i\} P\_\{i j\} = (P\textasciicircum{}2)\_\{u\_r j\} = P\_\{u\_r j\} = (p\_\{u\_r\})\_j. Using beta\_r(i) = P\_\{u\_r i\}, this gives p\_\{u\_r\} = sum\_i beta\_r(i) p\_i. \steptag{validated} \\[4pt]
\step{1.5.2} Substitute the defining chart expansion p\_i = sum\_q a\_q(i) p\_\{u\_q\} into p\_\{u\_r\} = sum\_i beta\_r(i) p\_i. The row index set and q in \{0,1,2\} are finite, so distributing scalars and interchanging the finite sums gives p\_\{u\_r\} = sum\_q (sum\_i beta\_r(i)*a\_q(i)) p\_\{u\_q\}. \steptag{validated} \\[4pt]
\step{1.5.3} The rows p\_\{u\_0\}, p\_\{u\_1\}, p\_\{u\_2\} form an ordered basis of the row space, so coordinates in that basis are unique. The standard expansion of p\_\{u\_r\} is p\_\{u\_r\} = sum\_q delta\_\{q r\} p\_\{u\_q\}. Comparing this expansion with p\_\{u\_r\} = sum\_q (sum\_i beta\_r(i)*a\_q(i)) p\_\{u\_q\} gives sum\_i beta\_r(i)*a\_q(i) = delta\_\{q r\} for every q in \{0,1,2\}, and in particular sum\_i beta\_r(i)*a\_s(i) = delta\_\{s r\}. \steptag{validated} \\[4pt]
\step{1.5.4} Because r and s are distinct, delta\_\{s r\} = 0. Therefore sum\_i beta\_r(i)*a\_s(i) = 0, with the sum over all row indices i of P, which is the asserted cross-pivot cancellation for this pair r,s. \steptag{validated} \\[4pt]
\step{1.5.5} Nodes 1.5.1 through 1.5.4 compose the local proof of node 1.5: 1.5.1 gives row reproduction for p\_\{u\_r\}; 1.5.2 converts that equality into chart coordinates; 1.5.3 identifies the s-coordinate as delta\_\{s r\} by basis uniqueness; and 1.5.4 uses r != s to make that coordinate zero. Therefore the statement of node 1.5 follows without using sibling node 1.4 as a dependency. \steptag{validated} \\[4pt]
\end{proofbox}

\end{document}
